\documentclass[11pt]{article}
\usepackage[utf8]{inputenc}
\usepackage[T1]{fontenc}
\usepackage{amsmath}
\usepackage{multicol}
\usepackage{geometry}
\usepackage{tikz}
\usetikzlibrary{shapes.geometric, arrows.meta}
\usepackage{enumitem}
\usepackage{xcolor}
\usepackage{titlesec}

% Configurações de layout
\geometry{a4paper, left=1cm, right=1cm, top=0.5cm, bottom=1.2cm}
\setlength{\columnseprule}{0.4pt}
\setlength{\baselineskip}{1.0\baselineskip}

% Cores para títulos
\titleformat{\section}{\normalfont\Large\bfseries\color{blue}}{\thesection}{1em}{}
\titleformat{\subsection}{\normalfont\large\bfseries\color{red}}{\thesubsection}{1em}{}
\titleformat{\subsubsection}{\normalfont\normalsize\bfseries\color{black}}{\thesubsubsection}{1em}{}

\title{\textcolor{blue}{Revisão: Cultura Digital e Contemporânea \\ Tema: Tecnologia e Sociedade \\ Análise de Alita - Anjo de Combate}}
\author{Professor: Jefferson}
\date{}

\begin{document}

\maketitle
\vspace{-1cm}

\begin{center}
\large{Nome: \underline{\hspace{8cm}} \quad Turma: \underline{\hspace{3cm}}}
\end{center}

\begin{multicols}{2}

\section*{Contexto}
O filme \textit{Alita: Anjo de Combate} (2019) apresenta um futuro distópico onde humanos e máquinas se fundem através da tecnologia ciborgue. Esta atividade propõe uma reflexão sobre como o filme aborda questões contemporâneas sobre tecnologia e sociedade.

\section*{Conceitos Fundamentais}

\subsection*{Ciborguização}
Fusão entre organismo biológico e tecnologia, representada por Alita, um ciborgue com cérebro humano e corpo mecânico.

\subsection*{Transumanismo}
Movimento que defende o uso da tecnologia para superar limitações humanas, tema central no filme.

\section*{Questões para Análise}

\begin{enumerate}[leftmargin=*, label=\arabic*., series=questoes]

\item Alita é um ciborgue com corpo mecânico e cérebro humano. Como essa representação se relaciona com os conceitos de ciborguização e transumanismo na cultura digital contemporânea?

\vspace{0.5cm}
\noindent\rule{\linewidth}{0.4pt}\\
\vspace{0.1cm}\\
\noindent\rule{\linewidth}{0.4pt}\\
\vspace{0.1cm}\\
\noindent\rule{\linewidth}{0.4pt}\\
\vspace{0.1cm}\\
\noindent\rule{\linewidth}{0.4pt}\\
\vspace{0.1cm}\\
\noindent\rule{\linewidth}{0.4pt}

\item A tecnologia apresentada no filme mostra a integração entre seres humanos e máquinas. Quais questões éticas e sociais esse fenômeno pode levantar em nossa realidade?

\vspace{0.5cm}
\noindent\rule{\linewidth}{0.4pt}\\
\vspace{0.1cm}\\
\noindent\rule{\linewidth}{0.4pt}\\
\vspace{0.1cm}\\
\noindent\rule{\linewidth}{0.4pt}\\
\vspace{0.1cm}\\
\noindent\rule{\linewidth}{0.4pt}\\
\vspace{0.1cm}\\
\noindent\rule{\linewidth}{0.4pt}

\item A memória fragmentada de Alita é um elemento central da narrativa. Como esse aspecto pode ser relacionado com nossos sistemas atuais de armazenamento e recuperação de dados digitais?

\vspace{0.5cm}
\noindent\rule{\linewidth}{0.4pt}\\
\vspace{0.1cm}\\
\noindent\rule{\linewidth}{0.4pt}\\
\vspace{0.1cm}\\
\noindent\rule{\linewidth}{0.4pt}\\
\vspace{0.1cm}\\
\noindent\rule{\linewidth}{0.4pt}\\
\vspace{0.1cm}\\
\noindent\rule{\linewidth}{0.4pt}

\item O esporte "Motorball" combina elementos físicos e tecnológicos. De que forma essa representação se conecta com as tendências atuais de gamificação e realidade aumentada?

\vspace{0.5cm}
\noindent\rule{\linewidth}{0.4pt}\\
\vspace{0.1cm}\\
\noindent\rule{\linewidth}{0.4pt}\\
\vspace{0.1cm}\\
\noindent\rule{\linewidth}{0.4pt}\\
\vspace{0.1cm}\\
\noindent\rule{\linewidth}{0.4pt}\\
\vspace{0.1cm}\\
\noindent\rule{\linewidth}{0.4pt}

\item Zalem representa uma cidade flutuante que controla tecnologicamente a cidade inferior. Que paralelos podemos estabelecer entre essa representação e as discussões atuais sobre concentração de poder tecnológico?

\vspace{0.5cm}
\noindent\rule{\linewidth}{0.4pt}\\
\vspace{0.1cm}\\
\noindent\rule{\linewidth}{0.4pt}\\
\vspace{0.1cm}\\
\noindent\rule{\linewidth}{0.4pt}\\
\vspace{0.1cm}\\
\noindent\rule{\linewidth}{0.4pt}\\
\vspace{0.1cm}\\
\noindent\rule{\linewidth}{0.4pt}

\item Os ciborgues do filme possuem partes mecânicas e humanas. Como essa representação dialoga com as pesquisas atuais em cibernética e interfaces cérebro-máquina?

\vspace{0.5cm}
\noindent\rule{\linewidth}{0.4pt}\\
\vspace{0.1cm}\\
\noindent\rule{\linewidth}{0.4pt}\\
\vspace{0.1cm}\\
\noindent\rule{\linewidth}{0.4pt}\\
\vspace{0.1cm}\\
\noindent\rule{\linewidth}{0.4pt}\\
\vspace{0.1cm}\\
\noindent\rule{\linewidth}{0.4pt}

\item A alta dependência tecnológica mostrada no filme apresenta quais riscos potenciais para a sociedade contemporânea?

\vspace{0.5cm}
\noindent\rule{\linewidth}{0.4pt}\\
\vspace{0.1cm}\\
\noindent\rule{\linewidth}{0.4pt}\\
\vspace{0.1cm}\\
\noindent\rule{\linewidth}{0.4pt}\\
\vspace{0.1cm}\\
\noindent\rule{\linewidth}{0.4pt}\\
\vspace{0.1cm}\\
\noindent\rule{\linewidth}{0.4pt}

\item A substituição de partes humanas por tecnologia, como mostrado no filme, traz benefícios médicos mas também dilemas. Discuta esses aspectos contraditórios.

\vspace{0.5cm}
\noindent\rule{\linewidth}{0.4pt}\\
\vspace{0.1cm}\\
\noindent\rule{\linewidth}{0.4pt}\\
\vspace{0.1cm}\\
\noindent\rule{\linewidth}{0.4pt}\\
\vspace{0.1cm}\\
\noindent\rule{\linewidth}{0.4pt}\\
\vspace{0.1cm}\\
\noindent\rule{\linewidth}{0.4pt}

\item O uso da tecnologia como forma de controle social no filme se assemelha a quais fenômenos digitais atuais?

\vspace{0.5cm}
\noindent\rule{\linewidth}{0.4pt}\\
\vspace{0.1cm}\\
\noindent\rule{\linewidth}{0.4pt}\\
\vspace{0.1cm}\\
\noindent\rule{\linewidth}{0.4pt}\\
\vspace{0.1cm}\\
\noindent\rule{\linewidth}{0.4pt}\\
\vspace{0.1cm}\\
\noindent\rule{\linewidth}{0.4pt}

\item O contraste entre a cidade tecnológica (Zalem) e a cidade inferior representa que tipo de divisão social no contexto da era digital?

\vspace{0.5cm}
\noindent\rule{\linewidth}{0.4pt}\\
\vspace{0.1cm}\\
\noindent\rule{\linewidth}{0.4pt}\\
\vspace{0.1cm}\\
\noindent\rule{\linewidth}{0.4pt}\\
\vspace{0.1cm}\\
\noindent\rule{\linewidth}{0.4pt}\\
\vspace{0.1cm}\\
\noindent\rule{\linewidth}{0.4pt}

\end{enumerate}

\section*{Atividade Prática}

\subsection*{Análise Comparativa}
Escolha uma tecnologia apresentada no filme e:
\begin{itemize}
\item Descreva seu funcionamento na narrativa
\item Identifique tecnologias reais equivalentes
\item Analise impactos sociais positivos e negativos
\item Proponha regulamentações éticas
\end{itemize}

\subsection*{Rubrica de Avaliação}
\begin{itemize}
\item Clareza da análise (0-3 pontos)
\item Relação com conceitos reais (0-3 pontos)
\item Profundidade da reflexão (0-3 pontos)
\item Originalidade das propostas (0-3 pontos)
\end{itemize}

\end{multicols}

\end{document}
